\documentclass[11pt]{article}
\usepackage{amsmath,amssymb,amsthm}
\usepackage{graphicx}
\usepackage{hyperref}
\usepackage{listings}

\newcommand{\R}{\mathbb{R}}
\newcommand{\norm}[1]{\left\lVert #1 \right\rVert}
\DeclareMathOperator{\tr}{tr}
\newtheorem{theorem}{Theorem}
\newtheorem{lemma}[theorem]{Lemma}

\title{A Short Paper on Sums}
\author{A. Writer \and B. Reader}
\date{\today}

\begin{document}
\maketitle

\begin{abstract}
We prove a closed form for the first $n$ whole numbers and note a bound in $\R^n$.
\end{abstract}

\section{Introduction}\label{sec:intro}

The sum of the first $n$ whole numbers was known to Gauss~\cite{gauss}.
It costs \$5 to print this paper, and 50\% of it is formulas.
We write \emph{emphasis}, \textbf{strong words} and \texttt{code} as usual,
and a footnote\footnote{A note at the foot, with $x^2$ in it.} too.
See Section~\ref{sec:result} for the result and \eqref{eq:gauss} for its statement.
More is at \url{https://example.org/sums} and in \href{https://example.org}{the archive}.

\section{The result}\label{sec:result}

\begin{theorem}[Gauss]\label{thm:gauss}
For every whole number $n \ge 1$,
\begin{equation}\label{eq:gauss}
  \sum_{i=1}^{n} i = \frac{n(n+1)}{2}.
\end{equation}
\end{theorem}

\begin{proof}
By induction. The case $n = 1$ is plain. If it holds for $n$, then
\begin{align}
  \sum_{i=1}^{n+1} i &= \frac{n(n+1)}{2} + (n+1) \\
                     &= \frac{(n+1)(n+2)}{2}. \nonumber
\end{align}
\end{proof}

\begin{lemma}\label{lem:norm}
For $x, y \in \R^n$ we have $\norm{x + y} \le \norm{x} + \norm{y}$, and
\[
  \tr\begin{pmatrix} a & b \\ c & d \end{pmatrix} = a + d .
\]
\end{lemma}

Theorem~\ref{thm:gauss} and Lemma~\ref{lem:norm} are all we need.

\subsection{Lists and a table}

\begin{itemize}
  \item A first point, with $\alpha \to \beta$.
  \item A second point:
  \begin{enumerate}
    \item numbered within;
    \item and again.
  \end{enumerate}
\end{itemize}

\begin{description}
  \item[Sum] what adding gives.
  \item[Product] what multiplying gives.
\end{description}

\begin{table}[h]
  \centering
  \begin{tabular}{l|rr}
    $n$ & sum & square \\ \hline
    1 & 1 & 1 \\
    2 & 3 & 4 \\
    3 & 6 & 9
  \end{tabular}
  \caption{The first three sums.}\label{tab:sums}
\end{table}

\begin{figure}[h]
  \centering
  \includegraphics[width=0.5\textwidth]{figures/triangle.png}
  \caption{A triangle of dots.}\label{fig:triangle}
\end{figure}

Table~\ref{tab:sums} and Figure~\ref{fig:triangle} show the pattern.

\begin{quote}
Mathematics is the queen of the sciences.
\end{quote}

\begin{verbatim}
def total(n):
    return n * (n + 1) // 2   # costs $0
\end{verbatim}

\begin{lstlisting}[language=Python]
print(total(100))
\end{lstlisting}

A drawing stands next, which Tome cannot show.

	\begin{tikzpicture}
	\draw (0,0) -- (1,1);
	\end{tikzpicture}

We also write \mycmd{x} within a line, and set x\quad y apart.
See \ref{nowhere}, the same book again \cite{gauss,missing}, and one page of it \cite[p.~3]{gauss}.

\section{Appendix}\label{sec:appendix}

This section is read from a second file, where the word zymurgy stands. It holds
one more formula, $\int_0^1 x\,dx = \tfrac12$, a command of its own \appcmd{y}, and a case split:
\[
  |x| = \begin{cases} x & \text{if } x \ge 0, \\ -x & \text{otherwise.} \end{cases}
\]


\section{Closing}\label{sec:closing}

This section is read from a file in a folder of its own. Section~\ref{sec:intro} began the paper.


\begin{thebibliography}{9}
\bibitem{gauss} C. F. Gauss, \emph{Disquisitiones Arithmeticae}, 1801.
\end{thebibliography}

\end{document}
